\chapter{Circuitos em série}\label{ch:g4:series-circuits}

No ano passado, um cordão de luzinhas de festa ficou apagado porque
uma única lâmpada faltava --- e prometemos que aquela arrumação tinha
um nome. Aqui está ele: o \omterm{def:g4:series-circuits:series}{circuito em série}, o jeito mais simples de
pôr muitos jogadores numa mesma volta. Ele tem regras rígidas da
casa, e uma lanterna na sua gaveta vem obedecendo a elas esse tempo
todo.

\section{Uma volta, vários jogadores}

\begin{definition}[Em série]\label{def:g4:series-circuits:series}
Os jogadores de um circuito estão \emph{em série}\index{circuito em
série} quando ficam numa única volta, um depois do outro, como contas
num mesmo colar: a estrada da corrente passa pelo primeiro, depois
pelo segundo, depois pelo terceiro, sem nenhuma bifurcação. Um
circuito montado assim é um \emph{circuito em série}.
\end{definition}

\begin{omfigure}
\begin{tikzpicture}[scale=0.9]
  \draw (0,0) to[battery1, l=pilha] (0,2.6)
    to[short] (1.6,2.6)
    to[cspst, l=interruptor] (3.2,2.6)
    to[lamp, l=lâmpada 1] (5.2,2.6)
    to[lamp, l=lâmpada 2] (7.2,2.6)
    to[short] (7.2,0)
    to[short] (0,0);
\end{tikzpicture}

{\small Um \omterm{def:g4:series-circuits:series}{circuito em série}: \omterm{def:g3:first-electric-circuit:battery}{pilha}, \omterm{ex:g3:first-electric-circuit:switch}{interruptor} e duas lâmpadas numa
única volta --- cada jogador enfiado no mesmo colar.}
\end{omfigure}

\begin{method}[Montar o circuito]\label{met:g4:series-circuits:build}
Com uma \omterm{def:g3:first-electric-circuit:battery}{pilha}, duas lâmpadas, um \omterm{ex:g3:first-electric-circuit:switch}{interruptor} e fios:
\begin{enumerate}
  \item comece no \omterm{def:g3:first-electric-circuit:battery}{terminal} $+$ da \omterm{def:g3:first-electric-circuit:battery}{pilha}, com um fio até o
        \omterm{ex:g3:first-electric-circuit:switch}{interruptor};
  \item leve um fio do \omterm{ex:g3:first-electric-circuit:switch}{interruptor} à primeira lâmpada, e da primeira
        lâmpada à segunda;
  \item feche a volta: da segunda lâmpada de volta ao \omterm{def:g3:first-electric-circuit:battery}{terminal} $-$ da
        \omterm{def:g3:first-electric-circuit:battery}{pilha};
  \item aperte o \omterm{ex:g3:first-electric-circuit:switch}{interruptor} e observe as duas lâmpadas juntas.
\end{enumerate}
Percorra a volta com o dedo antes de ligar: se em algum momento o seu
dedo tiver de escolher entre duas estradas, não é um \omterm{def:g4:series-circuits:series}{circuito em
série}.
\end{method}

\section{As regras da casa}

\begin{proposition}[Regras da volta única]\label{prop:g4:series-circuits:rules}
Num \omterm{def:g4:series-circuits:series}{circuito em série}:
\begin{enumerate}
  \item \emph{tudo ou nada}: uma falha em qualquer ponto --- um
        \omterm{ex:g3:first-electric-circuit:switch}{interruptor} aberto, uma lâmpada queimada, um fio solto --- e
        todos os jogadores param na hora;
  \item \emph{a ordem não importa}: troque as lâmpadas de lugar, mude
        o \omterm{ex:g3:first-electric-circuit:switch}{interruptor} para o outro lado da volta --- tudo se comporta
        exatamente como antes;
  \item \emph{repartição}: quanto mais lâmpadas na volta, mais fraco
        o brilho de cada uma; o empurrão da \omterm{def:g3:first-electric-circuit:battery}{pilha} é repartido pelo
        colar inteiro.
\end{enumerate}
\end{proposition}

\begin{example}[Vendo as regras]\label{ex:g4:series-circuits:seeing}
Monte o circuito de duas lâmpadas e teste cada regra. Desrosqueie
\emph{qualquer uma} das lâmpadas: as duas apagam --- regra um (o
soquete vazio é uma falha). Troque as duas lâmpadas de lugar, ou
ponha o \omterm{ex:g3:first-electric-circuit:switch}{interruptor} entre elas: nenhuma mudança --- regra dois. Agora
remonte com uma lâmpada só: ela brilha visivelmente mais forte
sozinha do que brilhava com a parceira --- regra três, ao contrário.
\end{example}

\begin{example}[As luzinhas de festa, resolvidas]\label{ex:g4:series-circuits:partylights}
O velho mistério agora deixou de ser mistério: as luzinhas de festa
estão \omterm{def:g4:series-circuits:series}{em série}, dezenas de contas num mesmo colar. Uma lâmpada
faltando é uma falha, e a regra um apaga o cordão inteiro. Encontrar
a lâmpada culpada testando uma por uma ensinou a regra um a gerações
de famílias, do jeito mais difícil --- os cordões modernos são feitos
de outro jeito, como o ano que vem vai mostrar.
\end{example}

\begin{remark}[O interruptor manda em todo mundo]\label{rem:g4:series-circuits:switch}
A regra dois tem uma consequência prática: um único \omterm{ex:g3:first-electric-circuit:switch}{interruptor},
colocado em qualquer ponto da volta, manda em todos os jogadores
dela. É por isso que um \omterm{ex:g3:first-electric-circuit:switch}{interruptor} de parede controla uma fileira
inteira de luminárias de teto --- e por isso, pelo mesmo motivo, você
não consegue apagar só \emph{uma} lâmpada de um cordão \omterm{def:g4:series-circuits:series}{em série} sem
apagar todas.
\end{remark}

\section{Mais pilhas no colar}

\begin{example}[Uma atrás da outra]\label{ex:g4:series-circuits:batteries}
As \omterm{def:g3:first-electric-circuit:battery}{pilhas} também podem entrar na fila \omterm{def:g4:series-circuits:series}{em série} --- mas elas são mais
exigentes que as lâmpadas: precisam ficar \emph{uma atrás da outra},
com o $+$ de cada \omterm{def:g3:first-electric-circuit:battery}{pilha} encostando no $-$ da seguinte. Duas \omterm{def:g3:first-electric-circuit:battery}{pilhas}
enfileiradas assim empurram juntas, e a lâmpada brilha nitidamente
mais forte do que com uma só. Vire uma das \omterm{def:g3:first-electric-circuit:battery}{pilhas} --- cara a cara ---
e os dois empurrões brigam entre si: a lâmpada não recebe nada.
\end{example}

\begin{omfigure}
\begin{tikzpicture}[scale=0.8]
  \draw (0,0) to[battery1] (0,1.8)
    to[battery1, l=duas pilhas] (0,3.6)
    to[short] (3.6,3.6)
    to[lamp, l=mais forte!] (3.6,0)
    to[short] (0,0);
  \begin{scope}[xshift=6.9cm]
    \draw (0,0) to[battery1, l=uma pilha] (0,1.8)
      to[short] (0,3.6)
      to[short] (3.6,3.6)
      to[lamp, l=mais fraca] (3.6,0)
      to[short] (0,0);
  \end{scope}
\end{tikzpicture}

{\small Duas \omterm{def:g3:first-electric-circuit:battery}{pilhas} enfileiradas uma atrás da outra empurram juntas:
a mesma lâmpada brilha mais forte do que com uma \omterm{def:g3:first-electric-circuit:battery}{pilha} sozinha.}
\end{omfigure}

\begin{example}[Dentro da lanterna]\label{ex:g4:series-circuits:flashlight}
Abra uma lanterna: as \omterm{def:g3:first-electric-circuit:battery}{pilhas} entram uma atrás da outra num tubo ---
uma fila pronta --- e depois a lâmpada e o \omterm{ex:g3:first-electric-circuit:switch}{interruptor} de botão
completam uma única volta pela carcaça. Uma lanterna \emph{é} um
\omterm{def:g4:series-circuits:series}{circuito em série} no bolso. E agora você sabe por que ela apaga se
uma única \omterm{def:g3:first-electric-circuit:battery}{pilha} for colocada ao contrário: cara a cara, os empurrões
se anulam.
\end{example}

\section{Exercícios}

\begin{exercise}[$\star$]\label{exo:g4:series-circuits:1}
O que quer dizer ``\omterm{def:g4:series-circuits:series}{em série}''? Qual imagem do dia a dia o capítulo
usa para isso?
\end{exercise}

\begin{exercise}[$\star$]\label{exo:g4:series-circuits:2}
Enuncie as três regras da casa do \omterm{def:g4:series-circuits:series}{circuito em série}.
\end{exercise}

\begin{exercise}[$\star$]\label{exo:g4:series-circuits:3}
No circuito de duas lâmpadas, a lâmpada 2 é desrosqueada. O que
acontece com a lâmpada 1, e por quê? Qual regra está em ação?
\end{exercise}

\begin{exercise}[$\star$]\label{exo:g4:series-circuits:4}
Zoe muda o \omterm{ex:g3:first-electric-circuit:switch}{interruptor} de antes da lâmpada 1 para depois da lâmpada
2. O que muda no comportamento do circuito? Qual regra diz isso?
\end{exercise}

\begin{exercise}[$\star$]\label{exo:g4:series-circuits:5}
Uma \omterm{def:g3:first-electric-circuit:battery}{pilha} acende bem uma lâmpada. Preveja o brilho com: duas lâmpadas
\omterm{def:g4:series-circuits:series}{em série}; \ três lâmpadas \omterm{def:g4:series-circuits:series}{em série}. O que está sendo repartido?
\end{exercise}

\begin{exercise}[$\star$]\label{exo:g4:series-circuits:6}
Como duas \omterm{def:g3:first-electric-circuit:battery}{pilhas} precisam ficar para empurrar juntas? Desenhe a fila
com o $+$ e o $-$ marcados em cada \omterm{def:g3:first-electric-circuit:battery}{pilha}.
\end{exercise}

\begin{exercise}[$\star$]\label{exo:g4:series-circuits:7}
Percorra a volta de uma lanterna: diga os jogadores que a corrente
visita do $+$ da \omterm{def:g3:first-electric-circuit:battery}{pilha} de volta até o $-$ dela.
\end{exercise}

\begin{exercise}[$\star$]\label{exo:g4:series-circuits:8}
Um circuito de campainha tem uma \omterm{def:g3:first-electric-circuit:battery}{pilha}, um botão e a campainha \omterm{def:g4:series-circuits:series}{em
série}. Onde você poderia acrescentar um segundo botão para que a
campainha toque só quando os \emph{dois} botões estiverem apertados?
\end{exercise}

\begin{exercise}[$\star$]\label{exo:g4:series-circuits:9}
Um cordão de dez luzinhas de festa \omterm{def:g4:series-circuits:series}{em série} está apagado. Todas as
lâmpadas parecem boas. Descreva um plano paciente, usando uma lâmpada
sabidamente boa, para encontrar a queimada.
\end{exercise}

\begin{exercise}[$\star\star$]\label{exo:g4:series-circuits:10}
A lanterna de Tom usa duas \omterm{def:g3:first-electric-circuit:battery}{pilhas} e brilha forte. Depois de trocar as
\omterm{def:g3:first-electric-circuit:battery}{pilhas} no escuro, ela não acende nada --- e mesmo assim as duas
\omterm{def:g3:first-electric-circuit:battery}{pilhas} novas estão cheias, a lâmpada está boa e nenhum fio está
solto. O que Tom fez, e por que a lanterna continua apagada? (Veja o
\cref{ex:g4:series-circuits:batteries}.)
\end{exercise}

\begin{exercise}[$\star\star$]\label{exo:g4:series-circuits:11}
Uma engenheira precisa colocar um botão de parada de emergência numa
máquina cujo motor fica numa única volta com a \omterm{def:g3:first-electric-circuit:battery}{pilha} dele. Ela diz:
``Qualquer ponto da volta serve.'' Ela está certa? E por que o plano
dela falharia se o circuito da máquina tivesse estradas com
bifurcação? (O capítulo do ano que vem tem o nome dessas estradas.)
\end{exercise}
